\chapter{INTRODUCCIÓN}
\label{cap:introduccion}

\section{Antecedentes}
\label{sec:antecedentes}

\subsection{Antecedentes generales}

La práctica regular de actividad física es una de las recomendaciones más
consistentes de la medicina preventiva: el \citet{acsm2021} recomienda que
los adultos combinen ejercicio cardiorrespiratorio con al menos dos sesiones
semanales de entrenamiento de fuerza que involucren los principales grupos
musculares. Los gimnasios se han convertido en el espacio más accesible para
cumplir esa recomendación, porque concentran en un mismo lugar máquinas de
resistencia guiada, poleas, bancos, racks y equipos de cardio.

Sin embargo, inscribirse en un gimnasio no garantiza la permanencia. En un
estudio de seguimiento de un año a nuevos miembros de clubes de
entrenamiento, \citet{gjestvang2020} analizaron los motivos y las barreras
que influyen en el inicio y el mantenimiento del ejercicio, y mostraron que
una parte importante de los usuarios deja de entrenar con regularidad en los
primeros meses. En la misma línea, \citet{sperandei2016} estudiaron los
registros de asistencia de miles de usuarios de un centro de
acondicionamiento físico en un entorno sin supervisión y encontraron que la
mayoría abandonaba la actividad antes de completar los primeros meses. En
conjunto, se estima que entre el 40\,\% y el 65\,\% de los usuarios nuevos
abandona el gimnasio durante los primeros meses, principalmente por la falta
de guía y de seguimiento personalizado.

La alternativa tradicional a ese problema es el entrenador personal, que
enseña la técnica, diseña la rutina y controla el progreso. Su costo, sin
embargo, deja fuera a la mayor parte de los usuarios, que terminan copiando
rutinas genéricas o usando las máquinas por imitación, con el riesgo de
lesión que eso implica.

Paralelamente, el teléfono inteligente se ha convertido en el dispositivo
personal por excelencia. En Bolivia, más del 90\,\% de los teléfonos móviles
utiliza el sistema operativo Android \citep{statcounter2025}, lo que convierte
a una aplicación Android en el canal con mayor alcance para llevar una guía de
entrenamiento al bolsillo de cada usuario. Las tecnologías actuales de
desarrollo multiplataforma, como React Native y Expo, y las plataformas de
servicios en la nube, como Firebase, permiten construir este tipo de
aplicaciones con un costo de infraestructura bajo y con tiempos de desarrollo
acotados.

\subsection{Antecedentes específicos}

En el mercado existen aplicaciones de entrenamiento que abordan parte del
problema. \emph{BodBot} \citep{bodbot} se presenta como un entrenador personal
basado en inteligencia artificial que genera planes de ejercicio a partir de
los objetivos y el equipamiento disponible del usuario; su modelo es de
suscripción y sus contenidos están en inglés. Las aplicaciones de
\emph{Leap Fitness Group} \citep{leapfitness}, muy difundidas en Android,
ofrecen rutinas predefinidas con animaciones, pero están orientadas al
entrenamiento en casa sin equipamiento, por lo que no enseñan el uso de las
máquinas de un gimnasio.

Ninguna de estas soluciones integra, en una misma aplicación en español y sin
costo para el usuario, los cuatro elementos que la literatura señala como
causas del abandono: la guía sobre el uso correcto de las máquinas, la
planificación estructurada del entrenamiento, el seguimiento del progreso
físico y una herramienta de control de la composición corporal. Esa
integración es el punto de partida del presente proyecto de grado.

\section{Definición del problema}
\label{sec:problema}

\subsection{Planteamiento del problema}

Los usuarios de gimnasio, en particular quienes recién comienzan, enfrentan
un conjunto de barreras que se refuerzan entre sí y que desembocan en el
abandono temprano de la actividad física:

\begin{itemize}
    \item \textbf{Desconocimiento del uso correcto de las máquinas y de la
    postura:} la ejecución por imitación produce errores técnicos que pueden
    causar lesiones y dificultan el progreso hacia los objetivos físicos.
    \item \textbf{Costo elevado de un entrenador personal:} la mayoría de los
    usuarios no puede pagar una guía profesional permanente.
    \item \textbf{Carencia de una planificación estructurada:} sin una rutina
    clara, el usuario no sabe qué grupos musculares trabajar cada día ni con
    cuántas series, repeticiones o tiempo de descanso, y desaprovecha su
    tiempo de entrenamiento.
    \item \textbf{Falta de seguimiento del progreso:} al no registrar su peso
    ni su composición corporal, el usuario no percibe sus avances, lo que
    reduce su motivación.
\end{itemize}

\subsection{Formulación del problema}

¿Cómo desarrollar una aplicación para gimnasios que ofrezca guía, rutinas de
ejercicio personalizadas mediante inteligencia artificial y seguimiento del
progreso físico, de forma accesible para el usuario?

\section{Objetivos}
\label{sec:objetivos}

\subsection{Objetivo general}

Desarrollar un sistema inteligente para el entrenamiento en gimnasios,
proporcionando una guía estructurada y seguimiento personalizado de los
ejercicios.

\subsection{Objetivos específicos}
\label{subsec:objetivos-especificos}

\begin{enumerate}[label=\textbf{OE\arabic*.}, leftmargin=3.2em]
    \item Guiar sobre el uso correcto de máquinas.
    \item Instruir en la ejecución correcta de ejercicios específicos.
    \item Diseñar rutinas de entrenamiento personalizadas y estructuradas,
    adaptadas a los objetivos de cada usuario.
    \item Construir herramientas de planificación que permitan a los usuarios
    organizar sus entrenamientos de manera eficiente.
    \item Registrar la ingesta de alimentos con el fin de controlar y evaluar
    las calorías consumidas por el usuario.
    \item Implementar un sistema de seguimiento del progreso físico.
    \item Generar reportes personalizados de progreso físico.
\end{enumerate}

\section{Área de conocimiento}
\label{sec:area}

\textbf{Área:} Programación Móvil.\\
\textbf{Subárea:} Sistemas de Información.

\section{Alcance y limitaciones}
\label{sec:alcance}

\subsection{Alcance}

\begin{itemize}
    \item El sistema funciona en dispositivos con sistema operativo Android.
    \item Registro e inicio de sesión de usuarios con correo electrónico y
    contraseña, o con una cuenta de Google, mediante una interfaz intuitiva.
    \item Guía sobre el uso correcto de las máquinas de gimnasio: un catálogo
    de 35 tipos de máquinas con descripción, músculos trabajados,
    instrucciones paso a paso, errores comunes, prevención de lesiones,
    galería de variantes y animaciones de ejecución de los ejercicios
    asociados.
    \item Escáner de máquinas con la cámara del teléfono, que superpone la
    ficha de la máquina sobre la imagen en vivo (realidad aumentada basada en
    cámara). La identificación de la máquina la confirma el usuario; la
    arquitectura deja definido el punto de extensión para un clasificador de
    imágenes.
    \item Creación de rutinas personalizadas según el objetivo y el nivel del
    usuario, generadas automáticamente por un sistema basado en reglas, o
    armadas manualmente día por día.
    \item Registro del progreso físico (peso, estatura e índice de masa
    corporal) con historial y reportes gráficos.
    \item El registro de alimentación se limita al control calórico mediante
    una calculadora de índice de masa corporal (IMC) con historial de
    mediciones.
\end{itemize}

\subsection{Limitaciones}

\begin{itemize}
    \item El sistema no reemplaza la supervisión profesional de entrenadores
    ni de nutricionistas.
    \item No se contempla la integración con dispositivos externos
    (relojes inteligentes, bandas de frecuencia cardíaca, entre otros).
    \item No se desarrolla una versión para iOS, aunque la tecnología elegida
    lo permitiría en el futuro.
    \item El reconocimiento automático de máquinas mediante un modelo de
    aprendizaje profundo (TensorFlow Lite) y la realidad aumentada con anclaje
    tridimensional (ARCore) quedan fuera del producto entregado; se documenta
    su diseño como trabajo futuro (véase la sección~\ref{sec:recomendaciones}).
    \item El contenido del catálogo de máquinas corresponde al equipamiento de
    un proveedor nacional de referencia y no a un gimnasio en particular.
\end{itemize}

\section{Justificación}
\label{sec:justificacion}

\subsection{Justificación técnica}

\begin{itemize}
    \item El uso de React Native con Expo y TypeScript permite obtener una
    aplicación Android nativa a partir de una sola base de código tipada, con
    acceso a la cámara y a los demás recursos del dispositivo.
    \item Firebase provee autenticación, base de datos documental en tiempo
    real y almacenamiento de archivos como servicios administrados, lo que
    elimina la necesidad de desplegar y mantener un servidor propio.
    \item La arquitectura por capas (Modelo--Vista--Controlador) aísla la
    lógica de negocio de la interfaz y de los servicios externos, de modo que
    componentes como el generador de rutinas o el reconocimiento de máquinas
    pueden evolucionar sin afectar al resto del sistema.
\end{itemize}

\subsection{Justificación económica}

\begin{itemize}
    \item La aplicación ofrece al usuario una guía de entrenamiento sin el
    costo recurrente de un entrenador personal.
    \item Las herramientas de desarrollo empleadas son libres o gratuitas, y
    los servicios de Firebase cuentan con un nivel gratuito suficiente para la
    etapa de validación.
    \item El catálogo se construyó a partir de fuentes de acceso libre (wger,
    free-exercise-db) y de contenido propio, sin licencias de pago.
\end{itemize}

\subsection{Justificación social}

\begin{itemize}
    \item Contribuye a reducir el abandono de la actividad física al atacar
    sus causas principales: desconocimiento, falta de estructura y falta de
    seguimiento.
    \item Disminuye el riesgo de lesiones al enseñar la ejecución correcta y
    los errores que deben evitarse en cada máquina.
    \item Pone a disposición de cualquier persona con un teléfono Android, y en
    español, una guía que antes dependía de poder pagar asesoría profesional.
\end{itemize}

\section{Metodología}
\label{sec:metodologia}

Para el desarrollo del producto se adoptó la metodología ágil Mobile-D
\citep{abrahamsson2004}, diseñada específicamente para proyectos de
aplicaciones móviles con equipos pequeños y ciclos de entrega cortos. El
trabajo se organizó en sus cinco fases: \emph{Exploración}, donde se
establecieron los requerimientos y se seleccionaron las tecnologías;
\emph{Inicialización}, donde se preparó el entorno, se definió la arquitectura
y se diseñó el modelo de datos; \emph{Producción}, donde se implementaron los
módulos funcionales en iteraciones sucesivas; \emph{Estabilización}, donde se
integraron los módulos y se refactorizó el código; y \emph{Pruebas del
sistema}, donde se verificó el producto contra los requerimientos. El
Capítulo~\ref{cap:desarrollo} describe cada fase en detalle.

En la parte documental se aplicó investigación bibliográfica sobre
prescripción del ejercicio, inteligencia artificial basada en conocimiento,
realidad aumentada y desarrollo móvil, que sustenta el marco teórico del
Capítulo~\ref{cap:marco}.
